
\examquestion{2023}{1}{A nonlinear optical mode: coherent states and Schr\"odinger cats}{2}
Consider the Hamiltonian
\[
H=\hbar[\omega a^\dagger a+\lambda(a^\dagger a)^p]
=\hbar[\omega\hat n+\lambda\hat n^p].
\]
Take \(p\) to be a \emph{positive} integer and
\(\lambda\ne0\) in the specified evolution time.

\subsection*{(a) Eigenstates and eigenvalues}
Because \(H\) is a function of \(\hat n=a^\dagger a\),
all number states are eigenstates:
\[
\boxed{H\ket n=E_n\ket n,\qquad
E_n=\hbar(\omega n+\lambda n^p),\quad n=0,1,2,\ldots}.
\]
This is a number-dependent nonlinearity. The
energy spacings are generally not constant.

\subsection*{(b) Evolve a coherent input in the interaction picture}
Take \(H_0=\hbar\omega\hat n\); it commutes with
\(H_I=\hbar\lambda\hat n^p\).
Write the initial coherent state as
\[
\ket\alpha=e^{-|\alpha|^2/2}
\sum_{n=0}^\infty\frac{\alpha^n}{\sqrt{n!}}\ket n .
\]
Since
\(U_I(t)=\exp(-i\lambda t\hat n^p)\),
the evolved state is
\[
\boxed{\ket{\psi_I(t)}=
e^{-|\alpha|^2/2}
\sum_{n=0}^\infty
\frac{\alpha^n}{\sqrt{n!}}
e^{-i\lambda t n^p}\ket n}.
\]
The nonlinear unitary changes phases but
\emph{not} photon-number probabilities.

\subsection*{(c) Recover the coherent-state cat for even \(p\)}
For \(p=2,4,\ldots\) and integer \(n\),
\(n^p\) is divisible by 4 if \(n\) is even,
while \(n^p\equiv1\pmod4\) if \(n\) is odd.
At the time \(t=\pi/(2\lambda)\),
\[
e^{-i\pi n^p/2}=
\begin{cases}1,&n\text{ even},\\-i,&n\text{ odd}.\end{cases}
\]
The parity-dependent coefficient is
\[
e^{-i\pi n^p/2}=
\tfrac12[(1-i)+(1+i)(-1)^n].
\]
Substitute into the Fock expansion and use
\((-1)^n\alpha^n=(-\alpha)^n\):
\[
\boxed{\ket{\psi_I}=\frac12
[(1-i)\ket\alpha+(1+i)\ket{-\alpha}]}.
\]
It is correctly normalised: the cross term
between \(\ket\alpha\) and \(\ket{-\alpha}\)
has purely imaginary coefficient, while
\(\braket{\alpha}{-\alpha}=e^{-2|\alpha|^2}\)
is real.

\subsection*{(d) Calculate all normal-ordered correlations}
The diagonal photon-number probabilities remain
\[
P_n=e^{-|\alpha|^2}\frac{|\alpha|^{2n}}{n!}.
\]
For positive integer \(k\), the normally
ordered factorial moment is
\[
\langle(a^\dagger)^ka^k\rangle
=\sum_{n=k}^\infty
n(n-1)\cdots(n-k+1)P_n=|\alpha|^{2k}.
\]
As \(\langle\hat n\rangle=|\alpha|^2\),
\[
\boxed{g^{(k)}(0)=
\frac{\langle(a^\dagger)^ka^k\rangle}{
\langle\hat n\rangle^k}=1
\qquad(k=1,2,\ldots,\ |\alpha|>0).}
\]
This is the usual \emph{normalised zero-delay}
intensity-correlation definition.

\subsection*{(e) Is the state classical?}
Despite having \emph{Poissonian photon statistics}
and \(g^{(k)}(0)=1\), the state is a coherent
superposition of distinguishable field amplitudes
\(\alpha\) and \(-\alpha\), not a classical mixture.
Its phase-space representation contains quantum
interference fringes and, for nonzero well-resolved
amplitudes, Wigner negativity; a positive classical
mixture of coherent states cannot reproduce it.
Photon counts alone fail to reveal these
phase-sensitive coherences. For \(\alpha=0\) the
state reduces to the vacuum (not a nontrivial cat).
\solcheck{The relative phase between Fock amplitudes
changes, but \(|e^{-i\lambda tn^p}|^2=1\),
so every photon-number probability is unchanged.}
